\subsection{Closed, closable, and densely defined operators}

In this section, K denotes a commutative topological field (TG, III, p. 54).

DEFINITION 2. --- Let E and F be topological vector spaces over K (EVT, I, p. 1, def. 1). Let \(u \in \mathcal{P}(\mathrm{E};\mathrm{F})\) be a partial operator from E to F.

We say that u is an operator with dense domain if the domain of u is dense in E.

We say that u is closed if the graph of u is closed in the topological vector space E × F. We say that u is closable if it has a closed extension.

Let E, F, and G be topological vector spaces over K. Every extension of an operator \(u \in \mathcal{P}(\mathrm{E};\mathrm{F})\) with dense domain has dense domain. Moreover, if \(v \in \mathcal{L}(\mathrm{E};\mathrm{F})\) , then \(u + v\) is an operator with dense domain. If \(v: F \to G\) (resp. \(w: G \to E\) ) is an isomorphism of topological vector spaces, then \(v \circ u\) (resp. \(u \circ w\) ) is an operator with dense domain.

Examples. --- Let E and F be topological vector spaces over K.

\begin{enumerate}
\def\labelenumi{\arabic{enumi})}
\item
  Assume the space F is Hausdorff. Let u be a linear map from E into F. If u is continuous, then the partial operator u is closed (TG, I, p. 53, cor. 2). Suppose moreover that K is a non-discrete valued field and that E and F are metrizable and complete topological vector spaces over K. By the closed graph theorem (EVT, I, p. 19, cor. 5), the partial operator defined by u is then closed if and only if u is continuous.
\item
  Suppose the space E is Hausdorff. Let v be an injective continuous linear map from F into E. The partial operator \(v^{-1} \in \mathcal{P}(\mathrm{E};\mathrm{F})\) (cf. IV, p. 226) is then closed, since its graph is the image of the graph of v, which is closed, by the isomorphism of topological vector spaces from \(F \times E\) into \(E \times F\) defined by \((y,x) \mapsto (x,y)\) .
\end{enumerate}

PROPOSITION 1. --- Let E and F be topological vector spaces over K. A partial operator u from E into F is closable if and only if the closure of the graph \(\Gamma_{u}\) of u in E × F is a functional graph. There then exists a unique partial operator v from E to F whose graph is \(\overline{\Gamma}_{u}\) , and it is the smallest closed extension of u.

If the closure of the graph of \(u\) in \(\mathrm{E} \times \mathrm{F}\) is a functional graph, it is the graph of a partial operator, and the latter is a closed extension of \(u\) , which is therefore closable. Conversely, suppose that \(u \subset w\) with \(w\) closed. The closure \(\overline{\Gamma}_u\) of the graph of \(u\) in \(\mathrm{E} \times \mathrm{F}\) is contained in \(\Gamma_w\) , hence \(\overline{\Gamma}_u\) is a functional graph.

The last assertion follows from the fact that if w is a closed extension of u, then the graph of w contains \(\overline{\Gamma}_{u}\) .

DEFINITION 3. --- Let E and F be topological vector spaces over K. Let u be a closable operator from E into F. The closed operator whose graph is \(\overline{\Gamma}_{u}\) is called the closure of \(u\) . It is denoted \(\overline{u}\) .

Remark. --- Let E and F be topological vector spaces over K. Let u be a closable operator from E to F. The domain of the closure of u is contained in the closure of the domain of u in E. In general, it is distinct from it (exercise 1 of IV, p. 344, b)).

If \(u \in \mathcal{P}(\mathrm{E};\mathrm{F})\) is closable and \(\operatorname{dom}(u) = \operatorname{E}\) , then \(u = \overline{u}\) is closed, since then \(\operatorname{dom}(\overline{u}) = \operatorname{dom}(u)\) .

PROPOSITION 2. --- Let K = R and let E and F be topological vector spaces over R. Let u be a partial operator from E into F. Then u is densely defined (resp. is closed, is closable) if and only if \(u_{(\mathbf{C})}\) has dense domain (resp. is closed, is closable).

Suppose that the domain of \(u\) is dense in E. Every neighborhood of 0 in \(\mathrm{E}_{(\mathbf{C})}\) contains a neighborhood of the form \(\mathrm{V} + i\mathrm{V}\) (EVT, II, p. 65), where V is a neighborhood of 0 in E, hence contains an element of the domain of \(u_{(\mathbf{C})}\) ; this partial operator is therefore densely defined. The converse is also true, since the map from \(\mathrm{E}_{(\mathbf{C})}\) into E that maps \(x + iy\) to \(x\) for all \((x,y) \in \mathrm{E} \times \mathrm{E}\) is continuous and surjective.

The graph of \(u_{(\mathbf{C})}\) is identified with the complexified topological vector space of the graph of \(u\) . It is therefore closed in \(\mathrm{E}_{(\mathbf{C})} \times \mathrm{E}_{(\mathbf{C})}\) if the graph of \(u\) is closed in \(\mathrm{E} \times \mathrm{E}\) . Conversely, we have \(\Gamma_u = \Gamma_{u_{(\mathbf{C})}} \cap (\mathrm{E} \times \mathrm{E})\) in \(\mathrm{E}_{(\mathbf{C})} \times \mathrm{E}_{(\mathbf{C})}\) ; since \(\mathrm{E} \times \mathrm{E}\) is closed in \(\mathrm{E}_{(\mathbf{C})} \times \mathrm{E}_{(\mathbf{C})}\) , the partial operator \(u\) is closed when \(u_{(\mathbf{C})}\) is.

A partial operator \(v\) from E into F is an extension of \(u\) if and only if \(v_{(\mathbf{C})}\) is an extension of \(u_{(\mathbf{C})}\) , so the partial operator \(u_{(\mathbf{C})}\) is closable if \(u\) is closable. Conversely, if \(u_{(\mathbf{C})}\) is closable, the partial operator \(u\) is also closable, since \(\overline{\Gamma_u} = \overline{\Gamma_{u_{(\mathbf{C})}}} \cap (\mathrm{E} \times \mathrm{E})\) , which is then a functional graph (prop. 1).

Lemma 1. — Let E, F, and G be topological vector spaces over K. Let u be a closed operator from E to F.

\begin{enumerate}
\def\labelenumi{\alph{enumi})}
\item
  For every \(v \in \mathcal{L}(\mathrm{E};\mathrm{F})\) , the partial operator \(u + v\) is closed.
\item
  For every \(v \in \mathcal{L}(\mathrm{G};\mathrm{E})\) , the partial operator \(u \circ v\) is closed.
\end{enumerate}

Let us prove \(a\) ). Let \(\gamma\) be the map \((x,y)\mapsto (x,y - v(x))\) from \(\mathrm{E}\times \mathrm{F}\) into itself; it is continuous. For every \((x,y)\in \mathrm{E}\times \mathrm{F}\) , we have \(\gamma (x,y)\in \Gamma_u\) if and only if \(x\in \operatorname {dom}(u)\) and \(y = u(x) + v(x)\) , that is, \(\overline{\gamma}^{1}(\Gamma_{u}) = \Gamma_{u + v}\) . The assertion follows.

Let us prove \(b\) ). The map \(\eta = (v, 1_{\mathrm{F}})\) from \(\mathbf{G} \times \mathbf{F}\) into \(\mathbf{E} \times \mathbf{F}\) is continuous; for every \((z, y) \in \mathbf{G} \times \mathbf{F}\) , we have \(\eta(z, y) = (v(z), y)\) , hence \(\overline{\eta}(\Gamma_u) = \Gamma_{u \circ v}\) . The assertion follows.

Let E and F be topological vector spaces over K, with F being Hausdorff. Let \(a \in \mathrm{K}\) . If \(u \in \mathcal{P}(\mathrm{E};\mathrm{F})\) is closable, then so is \(au\) . If \(a \neq 0\) , we have \(\overline{au} = a\overline{u}\) , and \(u\) is closed if and only if \(au\) is closed. If \(a = 0\) , the closure of \(au\) is \(0_{\overline{\operatorname{dom}(u)}}\) , and \(a\overline{u}\) is equal to \(0_{\operatorname{dom}(\overline{u})}\) ; it may therefore happen that \(u\) is closed but that \(au\) is not.

PROPOSITION 3. --- Let E and F be topological vector spaces over K, with F being Hausdorff. Let u be a closed partial operator from E to F. The kernel of u is a closed subspace of E.

Indeed, the kernel of u is the inverse image of the closed subspace \(\Gamma_{u} \cap (\mathrm{E} \times \{0\})\) of E × F under the continuous linear map \(x \mapsto (x, 0)\) from E into E × F.

In the rest of this section, we assume that K is a non-discrete valued field.

DEFINITION 4. --- Let E and F be normed spaces over K and let u be a partial operator from E to F. For x in dom(u), we denote

\[
\| x \| _ {u} = (\| x \| _ {\mathrm{E}} ^ {2} + \| u (x) \| _ {\mathrm{F}} ^ {2}) ^ {1 / 2}.
\]

The map \(x \mapsto \|x\|_{u}\) thus defined is a norm on \(\text{dom}(u)\) . We denote \(E_{u}\) the normed space thus obtained.

Remarks. --- Let E and F be normed spaces over K, and let u be a partial operator from E into F.

\begin{enumerate}
\def\labelenumi{\arabic{enumi})}
\item
  The canonical injection of \(E_{u}\) in E is continuous since we have \(\|x\|\leqslant\|x\|_{u}\) for all \(x\in E\) In particular, every subspace of \(\operatorname{dom}(u)\) that is closed in E is closed in \(E_{u}\) .
\item
  If E and F are Hilbert spaces, then the space \(E_{u}\) is a pre-Hilbert space, since the norm on \(E_{u}\) comes from the positive Hermitian form
\end{enumerate}

\[
(x, y) \mapsto (x \mid y) _ {u} = \langle x \mid y \rangle + \langle u (x) \mid u (y) \rangle
\]

on dom(u).

PROPOSITION 4. --- Let E and F be Banach spaces (resp. Hilbert spaces) and let u be a partial operator from E into F. Then u is closed if and only if the normed space \(E_{u}\) is a Banach space (resp. a Hilbert space).

It suffices to treat the case of Banach spaces. The norm of \(E_{u}\) is obtained, by transport of structure via the bijective linear map \((x,y)\mapsto x\) of \(\Gamma_{u}\) into dom(u), from the norm obtained by restriction to the subspace \(\Gamma_{u}\) with the norm \((x,y)\mapsto(\|x\|_{\mathrm{E}}^{2}+\|y\|_{\mathrm{F}}^{2})^{1/2}\) on the Banach space \(E\oplus F\) . The space \(E_{u}\) is therefore a Banach space if and only if the subspace \(\Gamma_{u}\) of \(E\oplus F\) is closed.

If \(u\) and \(v\) are partial operators from E to F and from F to G, respectively, and if \(u\) is closed, then the partial operator \(v \circ u\) is not closed in general, even if \(v\) is continuous (exercise 1 of IV, p. 344, c)). Nevertheless, we have the following sufficient condition:

Lemma 2. --- Let E, F, and G be normed spaces over K, with F being a Banach space. Let u be a closed partial operator from E to F, and let \(v \in \mathcal{L}(\mathrm{F};\mathrm{G})\) . If there exists \(\mathrm{C} \in \mathbf{R}_{+}\) such that

\[
\| u (x) \| \leqslant C (\| x \| + \| (v \circ u) (x) \|)
\]

for all \(x \in \text{dom}(v \circ u) = \text{dom}(u)\) , that is, if the linear map \(x \mapsto u(x)\) from \(E_{v \circ u}\) into F is continuous, then \(v \circ u\) is closed.

Denote by \(w = v \circ u\) . Let \((x_n, w(x_n))_{n \in \mathbf{N}}\) be a sequence in the graph of \(w\) which converges in \(E \times G\) . Let \(x\) be the limit of the sequence \((x_n)\) . The hypothesis implies that the sequence \((u(x_n))_{n \in \mathbf{N}}\) is then a Cauchy sequence in \(F\) ; it converges to an element \(y\) of \(F\) . Since \(u\) is closed, we have \(x \in \text{dom}(u)\) and \(y = u(x)\) . Then \(w(x_n) = v(u(x_n))\) tends to \(v(y) = v(u(x))\) since \(v\) is continuous, hence the graph of \(w\) is closed.

Let u be a closed partial operator on a Banach space E, and let F be a subspace of \(\text{dom}(u)\) . If F is dense in the Banach space \(E_{u}\) , then the reduction of u to F is closable, and its closure is equal to u. One then says that F is a core for u.

